\section{已有科研基础与所需的科研条件}

\subsection{已有科研基础}

本人围绕软件安全、漏洞分析与智能体安全方向开展了较为系统的文献调研和研究积累，已阅读并跟踪漏洞挖掘、自动漏洞修复、补丁验证、大语言模型与智能体应用于软件安全等领域的大量国内外研究工作，对传统程序分析方法、基于大模型的智能分析方法以及智能体驱动的软件安全研究范式具有较好的知识基础。通过前期调研，掌握了漏洞治理过程中漏洞发现、自动修复和修复验证等关键环节的发展现状与主要技术挑战，为后续开展基于智能体的漏洞挖掘与修复研究奠定了理论基础。

在漏洞修复验证方向，已开展了较为深入的研究工作。针对补丁存在性测试领域长期缺乏系统性评估、现有方法在真实部署场景下能力边界不明确的问题，完成了《A Comprehensive Empirical Analysis of Patch Presence Testing: Capabilities, Limitations, and Paths Forward》研究工作，并已被 \textit{ISSTA 2026} 录取~\cite{zhang2026pptempirical}。该工作首次针对补丁存在性测试开展系统性实证分析，构建大规模、高质量评测数据集，对现有代表性方法进行统一评估，分析影响检测效果的关键因素及失败原因，为后续智能体方法在补丁验证领域的发展提供了重要参考。

在上述研究基础上，进一步探索了基于智能体的端到端补丁存在性测试方法，完成了《Understanding Coding Agents for Patch Presence Testing》研究工作，目前已投稿~\cite{zhang2026agenticppt}。该工作将研究对象从传统专用分析工具扩展到通用编码智能体，系统分析智能体在补丁存在性测试任务中的能力边界，并提出面向该任务的智能体框架设计。实验结果表明，基于智能体的方法能够在真实二进制分析场景中达到与现有专用工具相当的性能，同时揭示了补丁定位和补丁演化识别等关键挑战，为后续研究智能体驱动的软件安全分析提供了实践基础。

在漏洞自动修复方向，已参与构建面向真实 Rust 漏洞的仓库级智能体修复评测基准《RRBench: A Repository-Level Benchmark for Agentic Repair of Real-World Rust Vulnerabilities》，目前已投稿~\cite{rrbench}。该工作针对现有漏洞自动修复评测中仅依赖 PoC 和功能测试、难以判断补丁是否真正恢复漏洞安全性质的问题，构建了包含真实漏洞案例、公开修复任务和独立安全验证机制的评测基准。通过对现有通用智能体修复能力的系统评估，发现当前智能体虽然能够生成大量表面有效的补丁，但仍存在根因理解不足、安全性质恢复不完整等问题。该工作为进一步研究基于智能体的漏洞自动修复方法、提升补丁生成可靠性提供了数据基础和问题依据。

综上，本人已在漏洞修复验证和智能体驱动漏洞分析方向形成了一定研究积累。其中，补丁存在性测试相关工作建立了对修复状态验证问题的系统认识，智能体补丁验证工作探索了端到端智能体解决方案，而漏洞修复评测基准工作进一步揭示了当前自动修复方法面临的关键挑战。这些研究基础为本课题后续围绕漏洞挖掘、漏洞自动修复以及漏洞修复验证三个方向开展深入研究提供了支撑。

\subsection{所需科研条件}

本课题研究需要的软件环境主要包括漏洞分析工具、程序分析框架、二进制分析工具以及大语言模型智能体开发环境。在漏洞挖掘方向，需要使用静态分析、动态分析、模糊测试和二进制分析工具，支持嵌入式设备、固件及复杂软件系统中的漏洞发现研究；在漏洞自动修复和修复验证方向，需要构建基于大语言模型的智能体框架，并结合代码分析、编译测试和漏洞验证环境开展实验。

硬件方面，需要具备支持大规模程序分析、大语言模型调用以及智能体实验运行的计算资源，包括高性能 CPU、GPU 服务器以及存储环境。同时，需要收集和维护真实开源项目、漏洞数据库、固件样本以及软件供应链组件等实验数据，为方法验证和系统评估提供基础。

目前，课题组已具备开展相关研究所需的软件安全研究环境，包括漏洞分析工具链、服务器计算资源以及大规模软件项目分析环境。同时，本人前期已积累漏洞数据构建、二进制分析、智能体框架开发和实验评测经验，能够支撑本课题后续研究工作的开展。
